\experiment{9}{%
    Develop a simple linear regression model for a dataset to establish the relationship between two variables. Use the model to make predictions and interpret the results.
}

\subsection{Building Linear regression model}

We'll training 2 linear regression model on our dataset:

1. x = Marks, y = StudyHours

2. x = Marks, y = Attendance

\subsubsection*{i) Model 1}

\begin{lstlisting}[style=python]
import pandas as pd 
from scipy import stats

m, b, r, p, se = stats.linregress(df_student['Marks'], df_student['StudyHours'])
print(f"Slope (m): {m:.4f}")
print(f"Intercept (b): {b:.4f}")
print(f"R^2 : {r**2:.4f}")
print(f"Equation: StudyHours = {m:.4f} * Marks + ({b:.4f})")
print(f"Standard Error: {se}")
\end{lstlisting}

\textbf{Output:}

Slope (m): 0.1059
Intercept (b): -3.3658
R² : 0.9701
Equation: StudyHours = 0.1059 * Marks + (-3.3658)
Standard Error: 0.003514849799273692

\gap

\begin{lstlisting}[style=python]
# Making Predictions
test_marks = [24, 28, 50, 60, 80, 90]
for mark in test_marks:
    predicted = m * mark + b
    print(f"Marks = {mark} -> Predicted StudyHours = {round(predicted, 2)}")
\end{lstlisting}

\textbf{Output:}

Marks = 24 → Predicted StudyHours = -0.82
Marks = 28 → Predicted StudyHours = -0.4
Marks = 50 → Predicted StudyHours = 1.93
Marks = 60 → Predicted StudyHours = 2.99
Marks = 80 → Predicted StudyHours = 5.1
Marks = 90 → Predicted StudyHours = 6.16


\subsubsection*{ii) Model 2}

\begin{lstlisting}[style=python]
m, b, r, p, se = stats.linregress(df_student['Marks'], df_student['Attendance'])

print(f"Slope (m): {m:.4f}")
print(f"Intercept (b): {b:.4f}")
print(f"R^2 : {r**2:.4f}")
print(f"Equation: Attendance = {m:.4f} * Marks + ({b:.4f})")
print(f"Standard Error: {se}")
\end{lstlisting}

\textbf{Output:}

Slope (m): 0.9319
Intercept (b): 10.1499
R² : 0.9963
Equation: Attendance = 0.9319 * Marks + (10.1499)
Standard Error: 0.010745191624759547

\gap

\begin{lstlisting}[style=python]
# Making Predictions
test_marks = [24, 28, 50, 60, 80, 90]
for mark in test_marks:
    predicted = m * mark + b
    print(f"Marks = {mark} -> Predicted Attendance = {round(predicted, 2)}\%")
\end{lstlisting}

\textbf{Output:}

Marks = 24 → Predicted Attendance = 32.51\%
Marks = 28 → Predicted Attendance = 36.24\%
Marks = 50 → Predicted Attendance = 56.74\%
Marks = 60 → Predicted Attendance = 66.06\%
Marks = 80 → Predicted Attendance = 84.7\%
Marks = 90 → Predicted Attendance = 94.02\%

\subsection{Observation}

\subsubsection*{Model 1: Marks vs StudyHours}

i. The regression equation is StudyHours = 0.1059 * Marks - 3.3658, it roughly indicate that for every 1 mark increase a student studying approximately 0.11 hours more.

ii. The R² value of 0.9701 indicates a very strong and reliable model, which is also confirmed by very low Standard Error.

iii. The model predicts negative study hour for marks below 32 due to lack of records in dataset.

iv. For marks within the dataset range (35--95), predictions are pretty reasonable.

\subsubsection*{Model 2: Marks vs Attendance}

i. The regression equation is Attendance = 0.9319 * Marks + 10.1499, meaning for every 1 mark increase corresponds to approximately 0.93\% increase in attendance.

ii. The R² value of 0.9963 is exceptionlly high, indicates this is the stronger of the two models, which is also confirmed by very low Standard Error.

iii. redictions within the data range are highly accurate

